%\vspace{-\acvSectionTopSkip}
\section{Berufserfahrung}

\cventry
{Rheinmetall Air Defence AG, Angestellt, on-site} % Job Title
{Software Engineer C++} % Organization
{Zürich, CH} % Location
{11/2024 - heute} % Date
\begin{cvitems}
\item{Entwickelte Core-Funktionalitäten der Battle-Management-Software (Oerlikon Skynex®) im 12-köpfigen Team – von Low-Level-Radarsteuerung bis UI-Architektur}
\item{Entwarf und implementierte taktische Kommunikationsprotokolle (TCP, REST, Protobuf) zur Echtzeit-Vernetzung von Radaren und Effektoren}
\item{Setzte komplexe QML-Oberflächen nach UX-Vorgaben um, die als Kernbedienoberfläche des Systems eingesetzt werden}
\item{Führte moderne Entwicklungsprozesse (Code Reviews, automatisierte Tests) für das gesamte Entwicklungsteam ein}
\item[]{\textbf{Frameworks/Tools: Modernes C++ (C++20), Qt, QML, REST, Lua, ant, Protobuf, TCP/IP, Boost, GTest, Enterprise Architect}}
\end{cvitems}

\cventry
{TRUMPF SE + Co. KG, Angestellt, on-site} % Job Title
{Software Engineer C++/C\#} % Organization
{Schramberg, DE} % Location
{03/2019 - 10/2024} % Date
\begin{cvitems}
\item{\textbf{Visionsystem:} Entwickelte ein C++-Visionsystem (Linux Realtime / Windows) maßgeblich weiter, realisierte eine Plugin-basierte Kameraplattform für OEM-Kunden mit Treibersoftware für Industriekameras und etablierte die CI/CD-Pipeline (Azure DevOps) sowie umfassende Testautomatisierung (Unit-, Integrations- und Benchmarktests)}
\item{\textbf{Quality Data Store:} Konzipierte und entwickelte ein .NET-basiertes System zur flexiblen Kundendaten-Ablage mit vielseitigen Datenbank-Integrationen, evaluierte und implementierte eine performante Client-Server-Kommunikation via gRPC und erstellte die vollumfängliche Architekturdokumentation (ISAQB)}
\item{\textbf{CAD/CAM-Prozesse:} Entwickelte ein CAD/CAM-Backend in C\# mit domänenspezifischer Sprache (DSL) und LionWeb im 20-köpfigen Entwicklungsteam}
\item{\textbf{Middleware \& Sensoranbindung:} Agierte als Scrum Master für ein 5-köpfiges Team, verantwortete die Sprint- sowie Kapazitätsplanung und entwickelte Features wie die OPC-UA/gRPC-Integration}
\item{\textbf{Qualitätssicherung:} Gewährleistete die durchgängige Softwarequalität über alle Projekte hinweg durch Pull-Request-Reviews, die Einhaltung von Clean-Code-Standards (SOLID, DRY), CI/CD-Pipelines (Azure DevOps) und statische Analyse (Sonarcloud)}
\item[]{\textbf{Frameworks/Tools: C++17, C\#, .NET, Python, Qt, Boost, Conan, gRPC, OPC-UA, OpenCV, GenICam, GTest, xUnit, CMake, Docker, Azure DevOps, Sonarcloud}}
\end{cvitems}
